\section{Ergebnisse}
\label{sec:ergebnisse}

\subsection{Experimentelles Setup}
In diesem Kapitel werden die Ergebnisse der in Kapitel~\ref{sec:methode} beschriebenen 
Experimente vorgestellt und analysiert. Untersucht wurde, inwiefern sich die Klassifikationsleistung
eines kompakten Student-Netzes durch Knowledge Distillation von einem größeren Teacher-Modell verbessern 
lässt. Als Vergleichsbasis dient das identische Student-Netz, das ohne Teacher ausschließlich mit den 
Ground-Truth-Labels trainiert wurde. Die Bewertung erfolgt anhand der Test- und Validierungsgenauigkeit, 
der Verlustverläufe während des Trainings sowie einer klassenweisen Analyse mittels Konfusionsmatrizen.

\subsection{Quantitative Ergebnisse}
\label{sec:quantitative_ergebnisse}
 
Tabelle~\ref{tab:distillation_ergebnisse} fasst die Ergebnisse aller Trainingsläufe zusammen. Insgesamt
liegen die erzielten Testgenauigkeiten in einem Bereich von 71{,}40\,\% bis 82{,}11\,\% und damit auf einem
grundsätzlich vergleichbaren Niveau. Innerhalb dieses Bereichs zeigen sich jedoch deutliche Unterschiede in
Abhängigkeit von der Wahl der Hyperparameter $\alpha$ (Gewichtung zwischen Distillations- und Klassifikationsverlust)
und $T$ (Temperatur der Softmax-Funktion).
 
Besonders hervorzuheben ist die Konfiguration $\alpha = 0{,}1$ und $T = 4{,}0$: Mit einer Testgenauigkeit von
82{,}11\,\% erzielt sie das mit Abstand beste Ergebnis aller Versuche. Gegenüber der Baseline von 73{,}16\,\% entspricht
dies einer absoluten Verbesserung von rund 9~Prozentpunkten bzw. einer relativen Steigerung von etwa 12\,\%. Dieses
Modell erreicht zudem mit 84{,}04\,\% die höchste Validierungsgenauigkeit und mit 0{,}671 den niedrigsten
Validierungsverlust aller Läufe, was auf eine konsistent bessere Generalisierung und nicht lediglich auf einen zufälligen
Ausreißer hindeutet.
 
Die beobachtete Verbesserung stellt sich jedoch nicht in allen Konfigurationen ein. Zwei Kombinationen --
$\alpha = 0{,}1$, $T = 2{,}0$ (71{,}40\,\%) sowie $\alpha = 0{,}3$, $T = 2{,}0$ (72{,}98\,\%) -- schneiden sogar
schlechter ab als die Baseline. Auffällig ist, dass beide unterlegenen Läufe mit der niedrigsten untersuchten
Temperatur $T = 2{,}0$ trainiert wurden, mit Ausnahme der Konfiguration $\alpha = 0{,}2$, $T = 2{,}0$, die mit
78{,}95\,\% überdurchschnittlich abschneidet. Weitere Konfigurationen, etwa $\alpha = 0{,}3$ mit $T = 6{,}0$
(74{,}39\,\%), liegen nur unwesentlich über der Baseline und weisen damit keine praktisch relevante Verbesserung auf.
Ein streng monotoner Zusammenhang zwischen den Hyperparametern und der erzielten Genauigkeit lässt sich aus den Daten
nicht ableiten; vielmehr deutet sich an, dass eine geringe Gewichtung des Hard-Label-Verlusts ($\alpha = 0{,}1$) in
Kombination mit einer moderaten Temperatur ($T = 4{,}0$) das Wissen des Teachers am effektivsten überträgt.
 
Hinsichtlich des Rechenaufwands ist festzuhalten, dass die Distillation die Trainingszeit von etwa 2{,}4~Minuten
auf 6{,}6 bis 7{,}7~Minuten nahezu verdreifacht, da in jedem Trainingsschritt zusätzlich ein Forward-Pass durch das
Teacher-Modell erforderlich ist. Da dieser Mehraufwand ausschließlich das Training betrifft und die Inferenzkosten des
Student-Netzes unverändert bleiben, ist er für den angestrebten Einsatzzweck eines kompakten Modells unkritisch.
 
Anzumerken ist außerdem, dass die absoluten Werte des Trainingsverlusts der destillierten Modelle
(vgl. Abbildungen~\ref{fig:train_loss} und \ref{fig:val_loss}) nicht direkt mit dem der Baseline vergleichbar sind, da sich die Verlustfunktion
der Distillation aus der gewichteten Summe von Kreuzentropie- und Kullback-Leibler-Term zusammensetzt und Letzterer
mit $T^2$ skaliert. Die durchweg höheren Trainingsverluste bei größeren Temperaturen sind daher konstruktionsbedingt
und kein Indiz für schlechtere Konvergenz. Die Validierungsverluste, die für alle Modelle einheitlich als Kreuzentropie
berechnet werden, bestätigen hingegen die Rangfolge der Genauigkeiten.
 
\subsection{Konfusionsmatrix und klassenweise Analyse}
 
Für eine differenziertere Betrachtung werden die Konfusionsmatrizen des besten destillierten Modells 
($\alpha = 0{,}1$, $T = 4{,}0$, Abbildung~\ref{fig:cm_distilled}) und der Baseline (Abbildung~\ref{fig:cm_baseline}) 
gegenübergestellt. Grundsätzlich weisen die Matrizen aller Läufe ein ähnliches Muster auf: Die Stärken und Schwächen 
konzentrieren sich auf dieselben Klassen und unterscheiden sich zwischen den Konfigurationen meist nur in Details.
 
Visuell eindeutig unterscheidbare Klassen wie \textit{beach}, \textit{cemetery}, \textit{chaparral}, \textit{freeway},
\textit{oil\_gas\_field}, \textit{oil\_well}, \textit{parking\_lot} und \textit{solar\_panel} werden von beiden Modellen
nahezu fehlerfrei klassifiziert. Die systematischen Verwechslungen treten hingegen erwartungsgemäß zwischen semantisch
und visuell verwandten Klassen auf. Hierzu zählen insbesondere:
 
\begin{itemize}
	\item \textit{sparse\_residential} und \textit{dense\_residential}, die sich primär in der Bebauungsdichte unterscheiden,
	\item \textit{basketball\_court} und \textit{tennis\_court}, die als rechteckige Sportflächen eine sehr ähnliche Struktur aufweisen,
	\item \textit{runway} und \textit{runway\_marking}, die denselben Bildinhalt in unterschiedlichem Maßstab zeigen,
	\item Klassen mit Wasser- bzw. Hafenbezug wie \textit{ferry\_terminal}, \textit{harbor} und \textit{bridge}.
\end{itemize}
 
Der Vergleich der beiden Matrizen zeigt, dass die Distillation vor allem bei diesen schwierigen Klassen zu Verbesserungen
führt. Besonders deutlich wird dies bei \textit{dense\_residential}: Während die Baseline hier lediglich 6 von 15
Testbildern korrekt zuordnet und die übrigen überwiegend auf benachbarte Bebauungsklassen verteilt, klassifiziert das
destillierte Modell alle 15 Bilder korrekt. Ähnliche Zugewinne zeigen sich bei \textit{nursing\_home}
(10 statt 6 korrekte Zuordnungen), \textit{parking\_space} (13 statt 6) und \textit{crosswalk} (15 statt 11). Diese
Beobachtung stützt die Hypothese, dass die weichen Wahrscheinlichkeitsverteilungen des Teachers gerade bei ähnlichen
Klassen zusätzliche Information über deren Beziehungen zueinander liefern, die aus den harten Labels allein nicht
hervorgeht.
 
Vereinzelt verschlechtert sich die Erkennung einzelner Klassen durch die Distillation geringfügig, etwa bei
\textit{storage\_tank} (11 statt 12 korrekte Zuordnungen) oder \textit{airplain}. In der Gesamtbilanz überwiegen 
die Verbesserungen jedoch klar, was sich in der um rund 9~Prozentpunkte höheren Gesamtgenauigkeit widerspiegelt.
 
\subsection{Interpretation und qualitative Beobachtungen}
 
Zusammenfassend lässt sich festhalten, dass die Wahl der Hyperparameter $\alpha$ und $T$ das Ergebnis der Knowledge
Distillation maßgeblich beeinflusst. Nicht jede Kombination führt zu der erhofften Verbesserung: Zwei der neun
untersuchten Konfigurationen blieben hinter der Baseline zurück, weitere erreichten lediglich ein vergleichbares Niveau.
In der Mehrzahl der Fälle bewirkte die Distillation jedoch eine Steigerung der Testgenauigkeit in der Größenordnung von
3 bis 5~Prozentpunkten; in der besten Konfiguration ($\alpha = 0{,}1$, $T = 4{,}0$) betrug der Zugewinn sogar
rund 9~Prozentpunkte.
 
Die Ergebnisse legen nahe, dass eine niedrige Gewichtung des Hard-Label-Anteils in Verbindung mit einer moderaten
Temperatur am besten geeignet ist, das im Teacher enthaltene Wissen über die Ähnlichkeitsstruktur der Klassen auf das
Student-Netz zu übertragen. Bei zu niedriger Temperatur nähern sich die weichen Zielverteilungen den harten Labels an,
sodass der zusätzliche Informationsgehalt weitgehend verloren geht; dies bietet eine plausible Erklärung für das schwache
Abschneiden mehrerer Läufe mit $T = 2{,}0$. Die klassenweise Analyse bestätigt darüber hinaus, dass der Nutzen der
Distillation vor allem bei visuell schwer unterscheidbaren Klassen zum Tragen kommt, während bei ohnehin gut separierbaren
Klassen kaum Spielraum für Verbesserungen besteht.
